\section{动量与自适应缩放改变的是动力学}
\label{sec:09-Convex-Optimization/08-Nonconvex-and-Stochastic-Optimization/05}

你在一个狭长的二次碗里跑 SGD。损失地形的等高线是极度拉长的椭圆：沿短轴方向曲率很大，沿长轴方向曲率很小。为了不让短轴方向上的更新过大而发散，你只能使用一个保守的学习率——而这使得长轴方向上的进展慢得难以忍受。每一步都在短轴方向上来回振荡，长轴方向缓慢推进，轨迹像锯齿。

你自然会想到：如果在持续下降的方向上积累势头，在来回震荡的方向上让正负梯度互相抵消，不就能在长轴方向走得更快吗？

这就是动量（momentum）。\hyperref[sec:14-Deep-Learning/03-Optimization-and-Training/02]{``动量与 Adam 在平均什么''}已经解释具体的更新式；本节关心这些附加状态如何改变优化的动力学——它们把一个一阶方法变成了更高阶的动力系统，收敛保证也因此不同。

\subsection{动量把一阶更新变成了二阶递推}

标准动量的未归一化写法：

\[
v_{t+1}=\beta v_t+\nabla f(x_t),
\qquad
x_{t+1}=x_t-\eta v_{t+1},
\]

其中 \(0\le\beta<1\) 控制旧速度的保留程度。消去速度变量，参数更新变成双步递推：

\[
x_{t+1}=(1+\beta)x_t-\beta x_{t-1}-\eta\nabla f(x_t).
\]

下一步不再只依赖当前参数和当前梯度——还依赖上一步的参数。这已经是一个带阻尼的离散动力系统，不是单纯的梯度下降。

在一维二次函数 \(f(x)=\frac{\lambda}{2}x^2\) 上可以精确求解这个动力系统的行为。代入 \(\nabla f(x_t)=\lambda x_t\)：

\[
x_{t+1}=(1+\beta-\eta\lambda)x_t-\beta x_{t-1}.
\]

设 \(x_t=r^t\)，代入得到特征方程：

\[
r^2-(1+\beta-\eta\lambda)r+\beta=0.
\]

稳定要求两个特征根都落在单位圆内。对 \(0\le\beta<1\)，稳定性条件为

\[
0<\eta\lambda<2(1+\beta).
\]

与普 SGD（\(\beta=0\)）的稳定条件 \(0<\eta\lambda<2\) 相比，动量把稳定区间从 2 扩展到了 \(2(1+\beta)\)——在学习率不变时，动量可以容忍更大的曲率。但反过来，对于固定的 \(\eta\lambda\)，过大的 \(\beta\) 会让特征根变成复共轭对，参数的迭代路径出现振荡——此时系统是稳定的，但收敛路径绕圈而不是单调下降。

动量在狭长碗中的收益来自各向异性：沿高曲率方向，有效 \(\eta\lambda\) 接近稳定上限，特征根接近单位圆——慢；沿低曲率方向，有效 \(\eta\lambda\) 小，特征根在圆内深处——快。动量主要加速了后者，对前者帮助有限。

需要警惕的是：不同实现中 \(1-\beta\) 被吸收到梯度还是学习率的方式不同（如 PyTorch 的 `dampening` 参数、不同框架的归一化惯例），上面给出的稳定范围只对应当前的记号。比较超参数前，先把各自的更新式统一到相同的形式。

\subsection{随机动量：滤波与记忆是一枚硬币的两面}

动量在随机梯度上的行为可以理解为指数移动平均滤波：梯度中频繁换号的高频分量被大量抵消，持续同号的低频分量被保留。这看起来像降噪——但有一个附加的成本：相邻 mini-batch 梯度的噪声结构不是``真实梯度加独立白噪声''。数据顺序引入了相关性；参数移动改变了每个样本的梯度；曲率矩阵本身也在旋转。

\(\beta\) 越大，有效记忆长度越长（约为 \(1/(1-\beta)\) 步），对持续方向的加速越强。但这也意味着，当目标函数的地形已经改变——比如越过了一个鞍点、进入了一个新的谷——旧的方向仍在影响更新。训练初期的 bias correction 和 warmup 策略，正是为了管理这个记忆在早期（速度尚未建立起可靠方向时的噪声放大）和后期（方向惯性导致对新地形的适应延迟）之间的权衡。

\subsection{自适应分母：从无偏梯度到有偏更新}

Adam 在动量之外添加了第二层状态：逐坐标维护梯度平方的指数移动平均。更新方向变成

\[
\frac{\widehat m_t}{\sqrt{\widehat v_t}+\varepsilon},
\]

其中 \(\widehat m_t\) 是梯度一阶矩的偏差修正估计，\(\widehat v_t\) 是二阶矩的偏差修正估计。

即使原始随机梯度本身满足条件无偏性 \(\E[g_t\mid\theta_t]=\nabla F(\theta_t)\)，这个比值一般也不是某个固定预条件矩阵乘真实梯度的无偏估计。分子和分母由同一段随机历史的样本耦合在一起——大梯度的坐标同时具有大的分子和大的分母，比值的行为不再是简单的一阶矩。

这对理论分析有直接影响：确定性梯度下降中依赖梯度无偏性的下降不等式，不能自动从 SGD 继承到 Adam。需要重新建立随机非凸设定下的一阶稳定性和收敛速率保证，而这些保证的成立条件比 SGD 的更严格。

Adam 还存在一个已知的问题：逐坐标有效步长可能因为旧的大梯度逐渐被指数平均遗忘而重新增大。在某些简单凸在线优化问题中，标准参数化下的 Adam 确实不收敛——算法会在最优解周围持续振荡而不进入更小的邻域。AMSGrad 用一个简单的修正应对这件事：始终取历史中最大的二阶矩

\[
\widetilde v_t=\max(\widetilde v_{t-1},v_t),
\]

使得有效步长只能单调不增，从而堵住了标准 Adam 证明链中的缺口，并在有界梯度和适当步长条件下恢复了某些 regret 上界。

但 AMSGrad 修复的是一个特定的理论缺陷，它不保证 Adam 或 AMSGrad 在所有非凸深度网络训练中都具有全局收敛保证——那样的保证目前不存在。它所做的是移除了一个已知的反例构造空间。

\subsection{经验优势与理论保证不可互相替代}

Adam 对稀疏梯度（如词嵌入中只有少数参数在每一步获得非零梯度）和跨坐标的尺度差异（不同层的梯度范数可能相差几个数量级）往往比 SGD 方便得多——调参区间更宽，默认参数在大量任务上有合理的表现。SGD with momentum 在某些视觉任务上可能收敛到不同的局部极小，泛化表现有时更好。这些都是反复观察到的经验规律。

但``在某个数据集上 loss 降得更快''不能等价于``该优化器具有更强的泛化保证''。更快的训练曲线可能是更快地收敛到了一个较差的极小值，或者优化器引入的隐式正则化与数据不匹配。报告优化器性能时，经验指标和理论属性应该分开陈述。

比较两个优化器时，以下因素中任何一个未控制都可能导致误导性结论：

\begin{itemize}
\tightlist
\item
  基础学习率与学习率 schedule（余弦退火、阶梯衰减、warmup 长度）；
\item
  batch size 与训练总步数（不同优化器的最优 batch size 可能不同）；
\item
  权重衰减的定义与实现——Adam 中 weight decay 是直接作用在参数上（AdamW）还是混入梯度平方平均（原版 Adam）会得出不同的隐式正则化；
\item
  数据遍历顺序与随机种子；
\item
  对超参数网格的覆盖密度——一个优化器在粗网格上被报告为``不敏感''可能只是因为其最优区域恰好在粗网格的覆盖范围外。
\end{itemize}

优化器是一种动力系统设计。理解其状态变量、稳定区域和噪声耦合方式，比对着一张 ``SGD vs Adam'' 的推荐表照搬超参数，更可能在新模型上做出正确的选择。
